\kbvsechead{(೨)\hspace{0.45em}ಗುಣವಾಚಕಗಳು :}

\vspace{2.5pt plus 3pt minus 0.5pt}\noindent\hspace*{1.2em}\kbvnum{(1)}\hspace{0.55em}ಮೈಸೂರಿನ ತಾತಯ್ಯ ಒಳ್ಳೆಯ ವ್ಯಕ್ತಿಯಾಗಿದ್ದರು.\par

\vspace{2.5pt plus 3pt minus 0.5pt}\noindent\hspace*{1.2em}\kbvnum{(2)}\hspace{0.55em}ಇದು ಹಳೆಯ ಮನೆ.\par

\vspace{2.5pt plus 3pt minus 0.5pt}\noindent\hspace*{1.2em}\kbvnum{(3)}\hspace{0.55em}ಕಪಿಗೆ ಉದ್ದವಾದ ಬಾಲವುಂಟು.\par

\vspace{2.5pt plus 3pt minus 0.5pt}\noindent\hspace*{1.2em}\kbvnum{(4)}\hspace{0.55em}ನನ್ನ ಮೇಜು ಚಿಕ್ಕದು.\par

\vspace{2.5pt plus 3pt minus 0.5pt}\noindent\hspace*{1.2em}\kbvnum{(5)}\hspace{0.55em}ದುಷ್ಟರಾದ ಜನರಿಂದ ದೂರವಿರಬೇಕು.\par

\indent ಮೇಲಿನ ವಾಕ್ಯಗಳಲ್ಲಿ ದಪ್ಪಕ್ಷರಗಳಲ್ಲಿರುವ ಪದಗಳನ್ನು ಗಮನಿಸಿ, ಒಳ್ಳೆಯ, ಹಳೆಯ, ಉದ್ದವಾದ, ಚಿಕ್ಕದು, ದುಷ್ಟರಾದ—ಇವು ಮನುಷ್ಯ, ಸ್ಥಳ, ಪ್ರಾಣಿ ಹಾಗೂ ವಸ್ತುಗಳ ಗುಣ, ರೀತಿ, ಸ್ವಭಾವಗಳನ್ನು ತಿಳಿಸುವ ; ಪದಗಳಾಗಿವೆ ; ಇವು ವಿಶೇಷಣಗಳು.\par

\indent “ಮೈಸೂರಿನ ತಾತಯ್ಯ ಒಳ್ಳೆಯ ವ್ಯಕ್ತಿಯಾಗಿದ್ದರು” ಎಂಬ ವಾಕ್ಯವನ್ನು ಆಧಾರವಾಗಿಟ್ಟು ಕೊಂಡು, “ಅವರು ಎಂತಹ ವ್ಯಕ್ತಿಯಾಗಿದ್ದರು?”, ಎಂಬುದಾಗಿ ಪ್ರಶ್ನೆ ಹಾಕಿದರೆ ಉತ್ತರ ರೂಪವಾಗಿ ಬರುವ ಪದ “ಒಳ್ಳೆಯ”, ಈ ಪದವು ತಾತಯ್ಯನವರ ಗುಣ ಅಥವಾ ಸ್ವಭಾವವನ್ನು ತಿಳಿಸುವುದು. ಇದೇ ಕ್ರಮವನ್ನು ಅನುಸರಿಸಿ, ಪ್ರತಿಯೊಂದು ವಾಕ್ಯದಲ್ಲೂ ನಾಮಪದದ ಮೇಲೆ “ಎಂತಹ” ಎಂಬ ಪ್ರಶ್ನೆ ಹಾಕಿ ಉತ್ತರ ಪಡೆಯಿರಿ ; ನಿಮಗೆ ಲಭಿಸುವ ಉತ್ತರವು ನಾಮಪದದ ವಿಶೇಷಣವೇ ಆಗಿರುವುದು.\par

\begin{kbvrule}
\noindent{\kbvsecfont ವ್ಯಾಖ್ಯೆ :}\hspace{0.5em}ಮನುಷ್ಯ, ಸ್ಥಳ, ಪ್ರಾಣಿ ಹಾಗೂ ವಸ್ತುಗಳ ಗುಣ, ರೀತಿ, ಸ್ವಭಾವಗಳನ್ನು ತಿಳಿಸುವ ವಿಶೇಷಣಗಳಿಗೆ ಗುಣವಾಚಕಗಳೆಂದು ಹೆಸರು.
\end{kbvrule}

\kbvsechead{(೩)\hspace{0.45em}ಸಂಖ್ಯಾವಾಚಕಗಳು :}

\indent ಒಂದು, ಎರಡು, ಮೂರು, ಹತ್ತು, ನೂರು ಇವೇ ಮೊದಲಾದ ನಾಮಪದಗಳು ಸಂಖ್ಯೆಗಳನ್ನು ಹೆಸರಿಸುತ್ತವೆ ; ಈ ನಾಮಪದಗಳಿಗೆ ಸಂಖ್ಯಾ ವಾಚಕಗಳೆಂದು ಹೆಸರು.\par

\vspace{2.5pt plus 3pt minus 0.5pt}\noindent\hspace*{1.2em}\kbvnum{(1)}\hspace{0.55em}ನನ್ನ ತೂಕ ಐವತ್ತು ಕಿಲೋಗ್ರಾಂ...\par

\vspace{2.5pt plus 3pt minus 0.5pt}\noindent\hspace*{1.2em}\kbvnum{(2)}\hspace{0.55em}ಈ ಪುಸ್ತಕದ ಬೆಲೆ ಎಂಟು ರೂಪಾಯಿಗಳು.\par
